\documentclass[10pt,a4paper]{amsart}
\usepackage[margin=19mm]{geometry}
\usepackage{amsmath,amssymb,mathtools}
\usepackage[scr=rsfs,cal=euler]{mathalfa}
\usepackage{xfrac}
\usepackage[unicode,hidelinks,bookmarks=false]{hyperref}
\newcommand{\ie}{i.e.\ }
\newcommand{\cf}{cf.\ }
\newcommand{\resp}{respectively\ }
\newcommand{\Subsection}{\subsection}
\providecommand{\nobreakdash}{\nobreak\hbox{-}}
\setlength{\emergencystretch}{2em}
\multlinegap0pt
\usepackage{bm}
\usepackage{bbm}
\usepackage{dsfont}
\usepackage{xspace}
\usepackage{cancel}
\usepackage{mathtools}
\usepackage{amsthm}
\makeatletter
\@ifundefined{theorem}{\newtheorem{theorem}{Theorem}}{}
\@ifundefined{proposition}{\newtheorem{proposition}[theorem]{Proposition}}{}
\makeatother
\DeclareMathOperator{\GL}{GL}
\DeclareMathOperator{\NN}{\mathbb{N}}
\DeclareMathOperator{\ZZ}{\mathbb{Z}}
\DeclareMathOperator{\cdeg}{cdeg}
\DeclareMathOperator{\fixed}{fixed}
\title[Color rules for cyclic wreath products and semigroup algebra]{Color rules for cyclic wreath products and semigroup algebras from projective toric varieties}
\author{Fabi\'an Levic\'an-Santib\'a\~nez, Marino Romero}
\date{Source: arXiv:2504.19008v2 (CC BY 4.0)}
\begin{document}
\maketitle

\noindent\emph{Source and licence.} Color rules for cyclic wreath products and semigroup algebras from projective toric varieties. Version arXiv:2504.19008v2 (CC BY 4.0), distributed under the Creative Commons Attribution 4.0 International licence, as recorded in the arXiv licence metadata for this exact version and shown on the abstract page https://arxiv.org/abs/2504.19008v2. Full rights notice: licenses/arxiv-2504.19008-CC-BY-4.0.txt.\par\smallskip

\noindent\textit{Editorial scope.} Counted unit: the Proposition whose source label is \texttt{proposition:polytope-character} in the article (source0.tex lines 1296--1303, source label \texttt{proposition:polytope-character}), delivered together with the article's own complete original proof of it (source0.tex lines 1305--1349, one \texttt{proof} environment of 45 lines). No mathematics is written, repaired or re-derived: statement and proof are the article's own text line for line. Required context retained uncounted: proposition:kpn-is-multigraded-zk-sn-module (lines 1248--1250). Every definition, display and label of those lines is retained unchanged. Omitted from this package and not redistributed: 1--1247, 1251--1295, 1304--1304, 1350--2544. Nothing inside a retained excerpt is omitted or altered: every retained body is delivered exactly as the article's lines are written, so no omission receipt of any kind accompanies this package. Citations: the retained excerpts carry no citation command at all, so no bibliography entry of the article is transcribed and no reference is redistributed. Reference adaptation: none was needed and none was made; every reference inside the delivered unit resolves inside it, so no unresolved reference or citation is produced. Printed slips kept literal: no printed word, symbol, hypothesis or number of the counted statement or of its proof was altered. Layout and class: the article's own class is \texttt{amsart} with packages including amsmath, amsthm, amssymb, amsfonts, amscd, bm, mathtools, float, graphicx, mdframed, tikz, verbatim, geometry, enumerate; the wrapper is the lane's pinned \texttt{amsart} editorial wrapper at 10pt on A4, which restates the theorem-like environments the retained text needs and adds the article's own macro definitions that the retained text needs (\texttt{\textbackslash GL}, \texttt{\textbackslash NN}, \texttt{\textbackslash ZZ}, \texttt{\textbackslash cdeg}, \texttt{\textbackslash fixed}), with extra wrapper packages (bm, bbm, dsfont, xspace, cancel, mathtools). The delivered numbering is the unit's own. Assets: no retained span contains a figure environment, a graphics-inclusion command, a PDF-page inclusion, a table or a TikZ picture; no image file is copied into this package, the package declares no asset dependency and no image file is redistributed with it. Licence: version arXiv:2504.19008v2 of this article is distributed under the Creative Commons Attribution 4.0 International licence, as recorded in the arXiv licence metadata for this exact version and shown on the abstract page https://arxiv.org/abs/2504.19008v2. A full rights notice is delivered at licenses/arxiv-2504.19008-CC-BY-4.0.txt. Characters: every retained excerpt is pure ASCII, so the delivered engine, XeLaTeX with its default Latin Modern text font, reports no missing character.
\begin{proposition}\label{proposition:kpn-is-multigraded-zk-sn-module}
    The affine semigroup algebra $K[P^{\times n}]$ admits the structure of a multigraded $Z_k \wr S_n$-module, where the gradings are the projective degree and the combinatorial degrees are given by a finite collection of weight vectors $\{\bm{a}_i = (\bm{a}_i')^{\times n}\}_{i \in [I]} \subseteq \ZZ^{mn}$.
\end{proposition}

\begin{proposition} \label{proposition:polytope-character}
    Consider the bigraded $Z_k \wr S_n$-module structure on $K[P^{\times n}]$ as in Proposition \ref{proposition:kpn-is-multigraded-zk-sn-module}. Then, for $\sigma = u_{a_1}\sigma_1 \cdots u_{a_n}\sigma_n \in Z_k \wr S_n$ of cycle type $\vec{\lambda}(\sigma) = \vec{\lambda}$, we have the following character:
    \[
    \chi^{K[P^{\times n}]}(\sigma) = \sum_{d \geq 0} t^d \prod_{i=0}^{k-1}
    \prod_{j = 1}^{\ell(\lambda^i)} \sum_{v \in dP \cap \mathbb{Z}^m} ~
    u_i^{|v|} q^{(\lambda^i_j) (\bm{a}' \cdot v)}.
    \]
\end{proposition}

\begin{proof}
    Fix $d \in \NN$. Let $K[P^{\times n}]_d$ be the $d$-th projective piece of $K[P^{\times n}]$. For $i = 1, \ldots, n$ and $v_i \in dP \cap \ZZ^m$, write
    \[
    x_{i, v_i} \coloneq \prod_{j = 1}^m x_{i, j}^{\pi_j(v_i)}.
    \]
    The set
    \[
    \mathcal{B}_d := \left \{ t^d \prod_{i = 1}^n x_{i, v_i}: \quad v_i \in dP \cap \ZZ^m, \text{ for all } i = 1, \ldots, n \right \}
    \]
    is the standard (monomial) basis of $K[P^{\times n}]_d$ (in particular, $|\mathcal{B}_d| = |dP^{\times n} \cap \ZZ^{mn}|$). Let $x = t^d \prod_{i = 1}^n x_{i, v_i} \in \mathcal{B}_d$. By definition,
    \begin{equation}\label{eq:image-basis-element}
        \sigma x = t^d \prod_{i = 1}^n u_{a_i}^{|v_i|} x_{\sigma_i, v_i},
    \end{equation}
    where $|v_i| = \sum_{j = 1}^m \pi_j(v_i)$. This gives a generalised permutation representation 
    \[
    \Pi_d: Z_k \wr S_n \to \GL(K\mathcal{B}_d).
    \]
    To calculate the coefficient of $t^d$ in the character $\chi^{K[P^{\times n}]}(\sigma)$ it suffices to count the fixed points $x_{\fixed}$ of the restriction $\Pi_d \downarrow_{S_n}^{Z_k \wr S_n}(\tau)$ with their corresponding roots of unity $u(x_{\fixed})$, where $\tau \in S_n$ is the permutation underlying $\sigma$. \\
    
    Let $D = \{(i, j) \in \ZZ^2: \quad 0 \leq i \leq k - 1, 1 \leq j \leq \ell(\lambda^i)\}$. It is not hard to see that these fixed points are the elements of $\mathcal{B}_d$ of the form
    \[
    x_{\fixed} = x_{\fixed}\left ((v_{ij})_{(i, j) \in D}\right ) = t^d \prod_{(i, j) \in D}\prod_{u_{a_k}\sigma_k \in c_j^i} x_{\sigma_{k}, v_{ij}}, \quad v_{ij} \in dP \cap \ZZ^m, \forall (i, j) \in D,
    \]
    where $c_j^i$ is the $C_i$-cycle of $\sigma$ indexed by $j$. In other words, to construct a fixed point $x_{\fixed}$, we must choose a point in $dP \cap \ZZ^m$ for each cycle $c_j^i$ of $\sigma$. \\
    
    By (\ref{eq:image-basis-element}), their corresponding roots of unity are
    \[
    u(x_{\fixed}) = \prod_{(i, j) \in D}\prod_{u_{a_k}\sigma_k \in c_j^i} u_{a_{\sigma_k}}^{|v_{ij}|} = \prod_{(i, j) \in D} u_i^{|v_{ij}|}.
    \]
    Their combinatorial degrees are
    \begin{align*}
        \cdeg(x_{\fixed}) 
        &= \cdeg \left ( \prod_{(i, j) \in D}\prod_{u_{a_k}\sigma_k \in c_j^i} x_{\sigma_{k}, v_{ij}} \right ) \\
        &= \sum_{(i, j) \in D} \sum_{u_{a_k}\sigma_k \in c_j^i} \cdeg(x_{\sigma_{k}, v_{ij}}) \\
        &= \sum_{(i, j) \in D} (\lambda_j^i) (\bm{a}' \cdot v_{ij}).
    \end{align*}
    Therefore,
    \begin{align*}
        [t^d]\chi^{K[P^{\times n}]}(\sigma) 
        &= \sum_{x_{\fixed}} u(x_{\fixed})q^{\cdeg(x_{\fixed})} \\
        &= \sum_{(v_{ij})_{i, j} \in (dP \cap \ZZ^m)^D} \prod_{(i, j) \in D} u_i^{|v_{ij}|}q^{(\lambda_j^i) (\bm{a}' \cdot v_{ij})} \\
        &= \prod_{(i, j) \in D} \sum_{v_{ij} \in dP \cap \ZZ^m} u_i^{|v_{ij}|}q^{(\lambda_j^i) (\bm{a}' \cdot v_{ij})}.
    \end{align*}
    The result follows.
\end{proof}

\end{document}
